\chapter{Finite monoid preliminaries}\label{ch:finitemonoid}

\begin{lemma}[One-sided inverses]\label{lem:finite-unit}
  \lean{mul_eq_one_comm, IsUnit.of_mul_eq_one_right}\leanok
  In a finite monoid, $ab = 1$ implies $ba = 1$; in particular either
  one-sided inverse relation makes both factors units. (Supplied by
  Mathlib; recorded because the cardinality bound
  Lemma~\ref{lem:localdiv-card} hinges on it.)
\end{lemma}
\begin{proof}
  $x \mapsto bx$ is injective ($a$ retracts it: $a(bx) = (ab)x = x$),
  hence surjective by finiteness; a preimage $q$ of $1$ gives $bq = 1$,
  and $q = (ab)q = a(bq) = a$, so $ba = 1$.
\end{proof}

\begin{lemma}[Idempotent power]\label{lem:idem-pow}
  \lean{KRTheory.exists_pow_idempotent}\leanok
  Every element $a$ of a finite monoid has a power $a^n$ with
  $n \ge 1$ and $a^n a^n = a^n$.
\end{lemma}
\begin{proof}
  Pigeonhole on $n \mapsto a^n$ gives $i < j$ with $a^i = a^j$; put
  $p := j - i \ge 1$. Then $a^{i+t+p} = a^{i+p} a^t = a^i a^t =
  a^{i+t}$ for every $t$, hence by induction
  $a^{i+t+mp} = a^{i+t}$ for every $m$. Take $n := p(i+1)$: then
  $n \ge i$, so writing $n = i + t_0$ we get
  $a^{2n} = a^{i + t_0 + (i+1)p} = a^{i+t_0} = a^n$.
\end{proof}

\begin{lemma}[Strict subtype count]\label{lem:card-subtype-lt}
  \lean{KRTheory.Nat.card_subtype_lt}\leanok
  A predicate on a finite type that fails at some point carves out a
  strictly smaller subtype.
\end{lemma}
\begin{proof}
  Transport of the corresponding \texttt{Fintype} fact along a
  classically chosen enumeration.
\end{proof}

\begin{remark}[Deferred to milestone 8]
  \notready
  The remaining \S 3.6 preliminary --- a monoid generated by units is a
  group --- is consumed only by the main induction and is deliberately
  not formalized yet.
\end{remark}
